\chapter{Reinforcement Learning Foundations}\label{ch:ml:reinforcement-learning-foundations}

An agent learns to sell a block of shares by trial and error in a simulator and beats the desk's schedule by 0.9 basis
points per lot. In the market the simulator was built from, it loses 0.6 basis points to the same schedule. The simulator
had one wrong detail, liquidity that alternated every period instead of persisting, and the agent had learned to exploit
it: it waited out every thin period for the deep one that the simulator always delivered next. \omterm{def:ml:reinforcement-learning-foundations:mdp}{Reinforcement learning} is
optimal control learned from experience instead of solved from a model (Book~4, chapter~9); what it learns is the
environment it was trained in. This chapter works on a problem small enough to solve exactly, so that every learner can be
measured against the truth: the value methods, the \omterm{def:ml:reinforcement-learning-foundations:mdp}{policy} methods, bandits, the evaluation of a new \omterm{def:ml:reinforcement-learning-foundations:mdp}{policy} from old logs,
and the gap between a simulator and a market. Chapter~18 applies them to execution and market making.

\section{Markov decision processes and the Bellman equations}

\begin{definition}[Reinforcement learning, Markov decision process, policy, reward, cumulative reward]\label{def:ml:reinforcement-learning-foundations:mdp}
\emph{Reinforcement learning}\index{reinforcement learning} learns how to act from the consequences of actions. A
\emph{Markov decision process}\index{Markov decision process} (MDP) is a set of states $x$, actions $u$, transition
probabilities $\P(x_{t+1}\mid x_t, u_t)$ and a \emph{reward}\index{reward} $g_{t+1}$ received after each action; the
\emph{cumulative reward}\index{cumulative reward} from time $t$ is $G_t = \sum_{k\ge0}\gamma^kg_{t+k+1}$, with a discount
$\gamma\in(0,1]$ ($\gamma = 1$ on the finite horizons of this chapter). A \emph{policy}\index{policy} $\pi(u\mid x)$ gives the
probability of each action in each state.
\end{definition}

\begin{definition}[Action-value function, Bellman equation]\label{def:ml:reinforcement-learning-foundations:bellman}
The \emph{action-value function}\index{action-value function} of a \omterm{def:ml:reinforcement-learning-foundations:mdp}{policy} is $Q^\pi(x, u) = \E_\pi[G_t\mid x_t = x, u_t = u]$,
the expected \omterm{def:ml:reinforcement-learning-foundations:mdp}{cumulative reward} of taking $u$ in $x$ and following $\pi$ afterwards; the value function of Book~4 is $V^\pi(x)
= \sum_u\pi(u\mid x)Q^\pi(x, u)$. The \emph{Bellman equation}\index{Bellman equation} is the recursion $Q^\pi(x, u) =
\E[g_{t+1} + \gamma V^\pi(x_{t+1})\mid x, u]$; for the optimal \omterm{def:ml:reinforcement-learning-foundations:mdp}{policy}, $Q^*(x, u) = \E[g_{t+1} + \gamma\max_{u'}Q^*(x_{t+1},
u')\mid x, u]$.
\end{definition}

The chapter's problem is to sell 10 lots in 10 periods. The state is the period, the lots left and the liquidity, deep or
thin, which persists from one period to the next with probability 0.8. Selling $u$ lots costs $\eta u^2$ basis points of
temporary impact, with $\eta = 0.5$ when liquidity is deep and 2 when it is thin; holding $q$ lots after a sale costs $0.2
q^2$ for the risk of the price moving; whatever is left is sold in the last period. The \omterm{def:ml:reinforcement-learning-foundations:mdp}{reward} is minus the cost. With 242
states, backward induction on the \omterm{def:ml:reinforcement-learning-foundations:bellman}{Bellman equation} (\cref{lst:ml:rl:solve}) solves it exactly: the optimal \omterm{def:ml:reinforcement-learning-foundations:mdp}{policy} costs
2.59 basis points per lot (risk included), selling 5 lots at once if liquidity is deep and 2 if it is thin. The desk's
schedule, the best one that ignores liquidity, costs 3.27; selling one lot a period costs 6.76, because of the risk of
holding the block.

\section{Value methods: temporal differences and Q-learning}

\begin{definition}[Temporal-difference learning, Q-learning]\label{def:ml:reinforcement-learning-foundations:td}
\emph{Temporal-difference learning}\index{temporal-difference learning} updates an estimate of a value towards a target built
from the next \omterm{def:ml:reinforcement-learning-foundations:mdp}{reward} and the current estimate of the next state's value, $V(x_t)\leftarrow V(x_t) + \eta_\alpha\delta_t$ with
$\delta_t = g_{t+1} + \gamma V(x_{t+1}) - V(x_t)$, without waiting for the end of the episode. \emph{Q-learning}\index{Q-learning}
applies it to action values with the target $g_{t+1} + \gamma\max_{u'}Q(x_{t+1}, u')$, which learns the optimal \omterm{def:ml:reinforcement-learning-foundations:mdp}{policy}'s values
while acting otherwise (Watkins and Dayan, 1992); SARSA uses the action actually taken next, and learns the values of the
\omterm{def:ml:reinforcement-learning-foundations:mdp}{policy} it follows.
\end{definition}

\Cref{lst:ml:rl:qlearning} is tabular \omterm{def:ml:reinforcement-learning-foundations:td}{Q-learning} with a step size of 0.1 and $\varepsilon$-greedy exploration (a random action
one time in ten). \Cref{fig:ml:rl:curves} shows the value gap of the learned \omterm{def:ml:reinforcement-learning-foundations:mdp}{policy}, measured exactly, against the number of
episodes, averaged over three seeds. After 1\,000 episodes \omterm{def:ml:reinforcement-learning-foundations:td}{Q-learning}'s greedy \omterm{def:ml:reinforcement-learning-foundations:mdp}{policy} is still 1.41 basis points per lot from
the optimum; after 10\,000 it is 0.044 away. SARSA, which learns the value of its own exploring \omterm{def:ml:reinforcement-learning-foundations:mdp}{policy}, is safer early
(0.28 at 1\,000 episodes) and ends at 0.051.

\section{Policy methods: gradients and actor-critic}

\begin{definition}[Policy gradient, actor--critic method]\label{def:ml:reinforcement-learning-foundations:pg}
A \emph{policy gradient}\index{policy gradient} method adjusts the parameters $\theta$ of a stochastic \omterm{def:ml:reinforcement-learning-foundations:mdp}{policy} in the direction
$\E[G_t\nabla_\theta\log\pi_\theta(u_t\mid x_t)]$, the gradient of expected \omterm{def:ml:reinforcement-learning-foundations:mdp}{cumulative reward} (REINFORCE; Williams, 1992); a
baseline subtracted from $G_t$ lowers the variance without biasing it. An \emph{actor--critic method}\index{actor--critic
method} replaces $G_t$ by the temporal-difference error of a learned value function (the critic), so that the \omterm{def:ml:reinforcement-learning-foundations:mdp}{policy} (the
actor) learns at every step.
\end{definition}

REINFORCE with a softmax \omterm{def:ml:reinforcement-learning-foundations:mdp}{policy} and a baseline per state learns slowly and noisily: 0.31 basis points from the optimum after
1\,000 episodes, 0.067 after 10\,000, and with a larger step size it settled on a worse \omterm{def:ml:reinforcement-learning-foundations:mdp}{policy} that it never left. The
actor-critic of \cref{lst:ml:rl:ac} learns fastest here (0.11 at 1\,000 episodes, 0.009 at 10\,000). \omterm{def:ml:reinforcement-learning-foundations:mdp}{Policy} methods matter
when actions are continuous or many, where a maximum over actions is expensive, and when a stochastic \omterm{def:ml:reinforcement-learning-foundations:mdp}{policy} is itself the
goal; deep versions replace the tables by networks (Mnih and co-authors, 2015, for value methods).

\begin{omfigure}
\begin{tikzpicture}
\begin{loglogaxis}[width=11cm,height=6cm,xlabel={training episodes},ylabel={gap to the optimum (bp per lot)},
  tick label style={font=\small},label style={font=\small},xmin=80,xmax=12500,ymin=0.005,ymax=3,grid=major,
  grid style={gray!25},legend columns=4,legend style={font=\footnotesize,at={(0.5,-0.25)},anchor=north,draw=none,
  fill=none,/tikz/every even column/.append style={column sep=6pt}},table/col sep=comma]
\addplot[omThm,very thick,mark=*] table[x=episodes,y=q]{figdata/ml/17-reinforcement-learning-foundations/curves.csv};
\addplot[omProp,very thick,mark=square*] table[x=episodes,y=sarsa]{figdata/ml/17-reinforcement-learning-foundations/curves.csv};
\addplot[gray,very thick,mark=triangle*] table[x=episodes,y=reinforce]{figdata/ml/17-reinforcement-learning-foundations/curves.csv};
\addplot[omDef,very thick,mark=diamond*] table[x=episodes,y=ac]{figdata/ml/17-reinforcement-learning-foundations/curves.csv};
\legend{Q-learning,SARSA,REINFORCE,actor-critic}
\end{loglogaxis}
\end{tikzpicture}

\omcaption{Value gap to the dynamic-programming optimum of the \omterm{def:ml:reinforcement-learning-foundations:mdp}{policy} learned after a number of episodes (the greedy \omterm{def:ml:reinforcement-learning-foundations:mdp}{policy}
for \omterm{def:ml:reinforcement-learning-foundations:td}{Q-learning} and SARSA, the softmax \omterm{def:ml:reinforcement-learning-foundations:mdp}{policy} for the others), mean of three seeds. Data:
\texttt{ml\_rl.learning\_curves}.}\label{fig:ml:rl:curves}
\end{omfigure}

\section{Bandits and exploration}

\begin{definition}[Exploration--exploitation trade-off, multi-armed bandit, contextual bandit]\label{def:ml:reinforcement-learning-foundations:bandit}
The \emph{exploration--exploitation trade-off}\index{exploration--exploitation trade-off} is the choice between the action that
looks best now and an action that teaches more. A \emph{multi-armed bandit}\index{multi-armed bandit} is the MDP with one
state: each round an arm is chosen and its \omterm{def:ml:reinforcement-learning-foundations:mdp}{reward} drawn; the loss against always playing the best arm is the regret. A
\emph{contextual bandit}\index{contextual bandit} observes a context (the order, the market) before choosing, and learns a
\omterm{def:ml:reinforcement-learning-foundations:mdp}{reward} model per arm; its actions do not change future contexts.
\end{definition}

Five execution algorithms save 0, 0.5, 1.0, 1.2 and 1.5 basis points per order on average, with 5 basis points of noise per
order. Over 10\,000 orders (\cref{fig:ml:rl:regret}), $\varepsilon$-greedy with $\varepsilon = 0.1$ loses 1\,727 basis points
of cumulative saving against always using the best algorithm, and uses it for 78\% of the last thousand orders; with
$\varepsilon = 0.01$ it explores too little and locks on the wrong algorithm in some runs (2\,026, 65\%). The upper
confidence bound rule, choosing the arm with the highest mean plus $c$ noise standard deviations times $\sqrt{\ln t/n_a}$
(Auer, Cesa-Bianchi and Fischer, 2002), loses 1\,325 with $c = \sqrt2$ and 944 with $c = 1$ (83\% and 87\%): it explores the
arms it is unsure of, not at random. Most of the loss is inevitable: an arm 0.3 basis points worse than the best, under 5 of
noise, needs about two thousand orders on each to be told apart at two standard errors.

\begin{omfigure}
\begin{tikzpicture}
\begin{axis}[width=11cm,height=6cm,xlabel={orders},ylabel={cumulative regret (bp)},tick label style={font=\small},
  label style={font=\small},xmin=0,xmax=10000,ymin=0,grid=major,grid style={gray!25},
  xtick={0,2000,4000,6000,8000,10000},xticklabel style={/pgf/number format/1000 sep={\,}},scaled x ticks=false,legend cell align=left,
  legend style={font=\footnotesize,at={(0.02,0.98)},anchor=north west,fill=white,draw=none},table/col sep=comma]
\addplot[omThm,very thick] table[x=t,y=eps1]{figdata/ml/17-reinforcement-learning-foundations/regret.csv};
\addplot[omThm!50,very thick,dashed] table[x=t,y=eps01]{figdata/ml/17-reinforcement-learning-foundations/regret.csv};
\addplot[omDef!60,very thick,dashed] table[x=t,y=ucb141]{figdata/ml/17-reinforcement-learning-foundations/regret.csv};
\addplot[omDef,very thick] table[x=t,y=ucb1]{figdata/ml/17-reinforcement-learning-foundations/regret.csv};
\legend{$\varepsilon$-greedy 0.1,$\varepsilon$-greedy 0.01,UCB $c = \sqrt2$,UCB $c = 1$}
\end{axis}
\end{tikzpicture}

\omcaption{Cumulative regret of four exploration rules choosing among five execution algorithms, mean of 20 runs. Data:
\texttt{ml\_rl.bandits}.}\label{fig:ml:rl:regret}
\end{omfigure}

\section{Off-policy evaluation and the simulator problem}

\begin{definition}[Off-policy evaluation, doubly robust estimator]\label{def:ml:reinforcement-learning-foundations:ope}
\emph{Off-policy evaluation}\index{off-policy evaluation} estimates the value of a target \omterm{def:ml:reinforcement-learning-foundations:mdp}{policy} from episodes generated by
another, the behaviour \omterm{def:ml:reinforcement-learning-foundations:mdp}{policy}, typically by reweighting each episode by the ratio of the two policies' probabilities of the
actions taken (importance sampling, Book~4, chapter~26). The \emph{doubly robust estimator}\index{doubly robust estimator}
adds to a model's estimate of the target's value the importance-weighted errors of the model on the logged \omterm{def:ml:reinforcement-learning-foundations:mdp}{rewards}; it is
unbiased if either the probabilities or the model are right, and has lower variance than importance sampling when the model
is close (Dudík, Langford and Li, 2011).
\end{definition}

A desk rarely gets to try a new \omterm{def:ml:reinforcement-learning-foundations:mdp}{policy} on real orders before it knows its value. Here the desk's logs come from its own
schedule with random deviations 30\% of the time (cost 3.80); the target is the optimal \omterm{def:ml:reinforcement-learning-foundations:mdp}{policy} with deviations 10\% of the
time, so that every action it can take was sometimes logged (true cost 2.77). \Cref{tab:ml:rl:ope} compares the estimators
on 100 sets of 300 logged episodes; the model given to the \omterm{def:ml:reinforcement-learning-foundations:ope}{doubly robust estimator} has the thin-liquidity impact half too
high and liquidity too changeable.

\begin{table}[htbp]
\centering\small
\begin{tabular}{@{}lccc@{}}
\toprule
estimator & bias & standard deviation & root mean square error\\
\midrule
per-decision importance sampling & 0.20 & 1.75 & 1.76\\
weighted importance sampling & 0.49 & 0.63 & 0.80\\
model alone & 0.30 & 0 & 0.30\\
doubly robust & $-0.03$ & 0.29 & 0.29\\
\bottomrule
\end{tabular}
\caption{Estimates of the target \omterm{def:ml:reinforcement-learning-foundations:mdp}{policy}'s cost (true value 2.77 bp per lot) from 300 logged episodes of the desk's
behaviour \omterm{def:ml:reinforcement-learning-foundations:mdp}{policy}, over 100 repetitions (bp per lot). Data: \texttt{ml\_rl.off\_policy}.}\label{tab:ml:rl:ope}
\end{table}

Importance sampling is unbiased and useless alone: products of ten probability ratios make a few episodes dominate. Weighting
them to sum to one trades variance for bias. The model alone is precise and wrong; the \omterm{def:ml:reinforcement-learning-foundations:ope}{doubly robust estimator} keeps the
model's precision and removes most of its bias (\cref{lst:ml:rl:dr}).

\begin{definition}[Sim-to-real gap]\label{def:ml:reinforcement-learning-foundations:simreal}
The \emph{sim-to-real gap}\index{sim-to-real gap} is the difference between a \omterm{def:ml:reinforcement-learning-foundations:mdp}{policy}'s performance in the simulator it was
trained in and in the real environment. It is largest for policies that exploit what the simulator gets wrong, which is
what an optimiser does with any error that pays.
\end{definition}

The hook's agent is \omterm{def:ml:reinforcement-learning-foundations:td}{Q-learning} trained for 10\,000 episodes in a simulator identical to the market except that liquidity
alternates every period, as it did in a calendar pattern in the data the simulator was calibrated on (Book~7, chapter~19). In
the simulator it costs 2.77 basis points per lot against the desk's 3.66. In the market it costs 3.83 against the desk's
3.27 and the optimum's 2.59: waiting for deep liquidity that does not come, it carries its risk longer and sells into thin
markets at the end. The desk's schedule, which assumed nothing about liquidity, lost nothing when the assumption was wrong.

\begin{method}[Before trusting a learned policy]\label{met:ml:rl:method}
\begin{enumerate}
\item Solve the smallest version of the problem exactly and check that the learner finds it.
\item Estimate the new \omterm{def:ml:reinforcement-learning-foundations:mdp}{policy}'s value from real logs by \omterm{def:ml:reinforcement-learning-foundations:ope}{off-policy evaluation}, with a \omterm{def:ml:reinforcement-learning-foundations:ope}{doubly robust estimator} and the
behaviour \omterm{def:ml:reinforcement-learning-foundations:mdp}{policy}'s logged probabilities; log them for every decision.
\item Evaluate the \omterm{def:ml:reinforcement-learning-foundations:mdp}{policy} in simulators that differ from the training one in the assumptions it might exploit (persistence,
impact, fills), and prefer policies whose value is stable across them.
\item Compare with the simplest schedule that assumes least, and deploy with limits (chapter~18).
\end{enumerate}
\end{method}

\section{Tutorial: learning to sell}\label{tut:ml:reinforcement-learning-foundations:tut}

\begin{tutorial}
\textbf{Goal.} Solve the liquidation MDP exactly, learn it four ways, run the bandits, evaluate a \omterm{def:ml:reinforcement-learning-foundations:mdp}{policy} from logs and train
in a wrong simulator. \textbf{End state:} \cref{fig:ml:rl:curves,fig:ml:rl:regret}, \cref{tab:ml:rl:ope}.
\begin{tutsteps}
\item \textbf{Backward induction.}
\omcode{code/firm/rlcore/firm_rlcore.py}{36}{46}{The Bellman optimality recursion on a finite horizon.}\label{lst:ml:rl:solve}
\item \textbf{\omterm{def:ml:reinforcement-learning-foundations:td}{Q-learning} and SARSA.}
\omcode{code/firm/rlcore/firm_rlcore.py}{128}{148}{Tabular Q-learning (and SARSA) with epsilon-greedy exploration.}\label{lst:ml:rl:qlearning}
\item \textbf{Actor-critic.}
\omcode{code/firm/rlcore/firm_rlcore.py}{210}{229}{A one-step actor-critic with a softmax policy.}\label{lst:ml:rl:ac}
\item \textbf{Doubly robust evaluation.}
\omcode{code/firm/rlcore/firm_rlcore.py}{290}{301}{The per-decision doubly robust estimator.}\label{lst:ml:rl:dr}
\item \textbf{Run} \texttt{ml\_rl.benchmarks()}, \texttt{learning\_curves()}, \texttt{bandits()}, \texttt{off\_policy()},
\texttt{sim\_to\_real()} and \texttt{fig\_rl.py}.
\end{tutsteps}
\textbf{What to change next.} Make the agent's own sales drain liquidity (the \texttt{drain} parameter) and retrain; give the
\omterm{def:ml:reinforcement-learning-foundations:ope}{doubly robust estimator} the true model and then a worse one.
\end{tutorial}

\section{Build: reinforcement-learning core}\label{bld:ml:reinforcement-learning-foundations:rlcore}

\begin{build}
\bfield{Purpose} \omterm{def:ml:reinforcement-learning-foundations:mdp}{Reinforcement learning} tested against exact answers, and policies evaluated before they trade.
\bfield{Interface} \texttt{FiniteMDP} with \texttt{solve}, \texttt{evaluate}, \texttt{q\_values};
\texttt{liquidation\_mdp(T, Q, eta, p\_stay, phi, drain)}; \texttt{MDPEnv} with \texttt{reset(seed)}, \texttt{step(u)};
\texttt{q\_learning}, \texttt{sarsa}, \texttt{greedy}, \texttt{reinforce}, \texttt{actor\_critic}, \texttt{softmax\_policy};
\texttt{bandit(means, sd, n, rule, seed, eps, c)}; \texttt{rollouts}, \texttt{ope\_is}, \texttt{ope\_wis}, \texttt{ope\_dr}.
\bfield{Rules} Every learner is seeded and reproducible; learned policies are evaluated exactly where the MDP is known;
logged data carry the behaviour \omterm{def:ml:reinforcement-learning-foundations:mdp}{policy}'s probabilities.
\bfield{Acceptance tests} \texttt{code/firm/rlcore/tests/}: backward induction matches brute-force enumeration on a tiny MDP;
exact evaluation matches Monte Carlo; \omterm{def:ml:reinforcement-learning-foundations:td}{Q-learning} reaches the optimum on a small problem; importance sampling is unbiased
and the \omterm{def:ml:reinforcement-learning-foundations:ope}{doubly robust estimator} has a lower spread, and is exact with a perfect model and deterministic transitions; UCB's regret grows more slowly than a random \omterm{def:ml:reinforcement-learning-foundations:mdp}{policy}'s.
\bfield{Stretch} Function approximation (a network for $Q$, chapter~18); LinUCB for \omterm{def:ml:reinforcement-learning-foundations:bandit}{contextual bandits}; \omterm{def:ml:reinforcement-learning-foundations:ope}{off-policy
evaluation} with estimated behaviour probabilities.
\end{build}

\begin{omsources}
\item R.~S.~Sutton and A.~G.~Barto, \emph{Reinforcement Learning: An Introduction}, 2nd ed., MIT Press, 2018.
\item C.~J.~C.~H.~Watkins and P.~Dayan, ``Q-learning'', \emph{Machine Learning} 8, 1992.
\item R.~J.~Williams, ``Simple statistical gradient-following algorithms for connectionist reinforcement learning'',
\emph{Machine Learning} 8, 1992.
\item P.~Auer, N.~Cesa-Bianchi and P.~Fischer, ``Finite-time analysis of the multiarmed bandit problem'', \emph{Machine
Learning} 47, 2002.
\item M.~Dudík, J.~Langford and L.~Li, ``Doubly robust policy evaluation and learning'', arXiv:1103.4601, 2011.
\item N.~Jiang and L.~Li, ``Doubly robust off-policy value evaluation for reinforcement learning'', arXiv:1511.03722, 2016.
\item V.~Mnih and co-authors, ``Human-level control through deep reinforcement learning'', \emph{Nature} 518, 2015.
\end{omsources}

\section{Exercises}

\begin{exercise}[$\star$]\label{exo:ml:reinforcement-learning-foundations:1}
In the last period the whole remainder must be sold. Write the optimal value of the state (period 9, 3 lots, thin
liquidity) and of (period 9, 3 lots, deep).
\end{exercise}

\begin{exercise}[$\star$]\label{exo:ml:reinforcement-learning-foundations:2}
A \omterm{def:ml:reinforcement-learning-foundations:td}{Q-learning} update has $Q(x, u) = -10$, \omterm{def:ml:reinforcement-learning-foundations:mdp}{reward} $-3$, $\max_{u'}Q(x', u') = -5$ and step size 0.1. What is the new $Q(x,
u)$? What would SARSA use if the next action taken had $Q(x', u') = -8$?
\end{exercise}

\begin{exercise}[$\star$]\label{exo:ml:reinforcement-learning-foundations:3}
Why does the cost of selling one lot a period (6.76) exceed the desk's schedule (3.27) so much in this problem?
\end{exercise}

\begin{exercise}[$\star\star$]\label{exo:ml:reinforcement-learning-foundations:4}
Why is SARSA ahead of \omterm{def:ml:reinforcement-learning-foundations:td}{Q-learning} after 1\,000 episodes and behind after 10\,000?
\end{exercise}

\begin{exercise}[$\star\star$]\label{exo:ml:reinforcement-learning-foundations:5}
Why does ordinary importance sampling have such a large standard deviation here? Estimate the typical size of the weight of
an episode.
\end{exercise}

\begin{exercise}[$\star\star$]\label{exo:ml:reinforcement-learning-foundations:6}
\emph{Find the flaw.} ``Our RL execution agent beat TWAP by 2 basis points over a year of simulation, with a $t$-statistic of
40, so we are deploying it.''
\end{exercise}

\begin{exercise}[$\star\star\star$]\label{exo:ml:reinforcement-learning-foundations:7}
\emph{Coding.} Set \texttt{drain = 0.1} in the true market (each lot sold raises the chance that the next period is thin by
0.1) and compare the optimal cost and first action with those of the chapter's market. Explain the difference.
\end{exercise}

\begin{exercise}[$\star\star\star$]\label{exo:ml:reinforcement-learning-foundations:8}
Show that the \omterm{def:ml:reinforcement-learning-foundations:ope}{doubly robust estimator} is unbiased when the model is exact ($\hat Q = Q^{\pi_e}$), whatever the importance
weights, for a one-step problem.
\end{exercise}

\section{Problem: Learning to Sell}

\begin{problem}[{Weekend problem --- a policy is only as good as its world}]\label{pb:ml:reinforcement-learning-foundations:1}
The chapter's liquidation problem, learners, bandits, logs and simulator.

\medskip\noindent\textbf{Part I --- The exact problem.}
\begin{enumerate}
\item Describe the states, actions, \omterm{def:ml:reinforcement-learning-foundations:mdp}{rewards} and transitions.
\item What does the optimal \omterm{def:ml:reinforcement-learning-foundations:mdp}{policy} cost, and what does it do first in deep and in thin liquidity?
\item What do the desk's schedule and one lot a period cost?
\item Why is backward induction possible here and not in a real market?
\end{enumerate}

\medskip\noindent\textbf{Part II --- Learning it.}
\begin{enumerate}[resume]
\item What gap to the optimum does each learner leave after 1\,000 and 10\,000 episodes?
\item Why can REINFORCE settle on a worse \omterm{def:ml:reinforcement-learning-foundations:mdp}{policy} with a larger step size?
\item What does the actor-critic gain over REINFORCE?
\item How many episodes is 10\,000 in trading days, if an episode is one parent order?
\end{enumerate}

\medskip\noindent\textbf{Part III --- Bandits and logs.}
\begin{enumerate}[resume]
\item What regret does each exploration rule accumulate, and which is best?
\item Why does $\varepsilon = 0.01$ do worse than 0.1?
\item Compare the four \omterm{def:ml:reinforcement-learning-foundations:mdp}{off-policy} estimates.
\item What must the desk log for \omterm{def:ml:reinforcement-learning-foundations:ope}{off-policy evaluation} to be possible?
\end{enumerate}

\medskip\noindent\textbf{Part IV --- The verdict.}
\begin{enumerate}[resume]
\item State the \emph{named result}: \omterm{def:ml:reinforcement-learning-foundations:td}{Q-learning}'s value gap to the dynamic-programming optimum after 10\,000 episodes, and the
bias and spread of the three \omterm{def:ml:reinforcement-learning-foundations:mdp}{off-policy} estimators.
\item What did the agent trained in the wrong simulator learn, and what did it cost?
\item How would you have caught it before deployment?
\item When would you prefer a bandit to full \omterm{def:ml:reinforcement-learning-foundations:mdp}{reinforcement learning} for an execution choice?
\item Why is the desk's schedule robust?
\item What changes with continuous prices and sizes?
\item What role should a simulator play in developing a trading \omterm{def:ml:reinforcement-learning-foundations:mdp}{policy}?
\item In one sentence: what does a reinforcement-learning agent learn?
\end{enumerate}
\end{problem}

\section{Interview questions}

\begin{interviewq}[$\star$ \iqroles{researcher, mle}]\label{iq:ml:reinforcement-learning-foundations:1}
Write the Bellman optimality equation for action values and explain each term.
\end{interviewq}

\begin{interviewq}[$\star\star$ \iqroles{mle}]\label{iq:ml:reinforcement-learning-foundations:2}
What is the difference between \omterm{def:ml:reinforcement-learning-foundations:td}{Q-learning} and SARSA? When does it matter?
\end{interviewq}

\begin{interviewq}[$\star\star$ \iqroles{researcher}]\label{iq:ml:reinforcement-learning-foundations:3}
How would you estimate the value of a new execution \omterm{def:ml:reinforcement-learning-foundations:mdp}{policy} from historical fills, without running it?
\end{interviewq}

\begin{interviewq}[$\star\star$ \iqroles{researcher, trader}]\label{iq:ml:reinforcement-learning-foundations:4}
You must choose among five brokers' algorithms for each order. Design the choice rule.
\end{interviewq}

\begin{interviewq}[$\star\star$ \iqroles{mle}]\label{iq:ml:reinforcement-learning-foundations:5}
Explain the \omterm{def:ml:reinforcement-learning-foundations:pg}{policy gradient} theorem and why a baseline helps.
\end{interviewq}

\begin{interviewq}[$\star\star\star$ \iqroles{researcher}]\label{iq:ml:reinforcement-learning-foundations:6}
An agent trained in your market simulator beats every benchmark. List what could make it fail live, and how you would test
each.
\end{interviewq}
